\section{The open steps}\label{sec:ch25:open}

By Theorem~\ref{thm:poisson} (\cite[Theorem 3.7 and Remark 3.10]{P26}), step 7 is exactly a zero free half plane $\re s>\alpha$ with $\alpha<1$, a quasi Riemann hypothesis. Until 2026 none was known. In October 2026 OpenAI released a proof, produced with one of its AI models, that $\zeta(s)\ne0$ for $\re s>\frac78$, and the same for every Dirichlet $L$ function \cite{OAI26a}, together with a formal proof of both statements in the Lean proof assistant and a separate proof of the weaker half plane $\re s>\frac{11}{12}$ \cite{OAI26b}. The argument proves a power saving for smoothed sums of the M\"obius function twisted by Hecke characters of $\mathbb Q(\sqrt{-3})$ and uses nothing from the Thorin picture. It has not yet been refereed. We rebuilt the formal proof from the public repository (commit \texttt{adc7f12}, Lean 4.34.1, Mathlib \texttt{d13f23b}): the Lean kernel accepts both proofs, which use only the axioms \texttt{propext}, \texttt{Classical.choice} and \texttt{Quot.sound}, and the comparator of the Lean Focused Research Organization, run without its sandbox, confirms that the theorems proved are the published challenge statements. If it stands, step 7 holds with $\alpha=\frac78$, and the constructions of this book that follow from it, conditional on that unrefereed result, are listed in Table~\ref{tx5:tab:map}, Theorem~\ref{pol:thm:qrh} and Corollary~\ref{qrh:cor:thorin} of Volume II. Below $\frac78$ nothing changes. Positivity of the Thorin density on the single line $\re s=\alpha$ is a weaker statement than a zero free half plane (Remark~\ref{rem:ch2:oneline}). So the first theorem the Thorin path needs beyond the announcement, a zero free half plane $\re s>\alpha$ with $\alpha<\frac78$, is an open problem in its own right; the L\'evy form, the scale mixture, completeness and Riesz's criterion (Sections~\ref{ch:basepoint}, \ref{ch:strip} and \ref{ch:riesz}) restate it and do not get round it. The Thorin formulation offers two tools for it. The sine squared form, compensated as in Proposition~\ref{bp:prop:below}, writes the required positivity as domination of the pole piece by the prime and archimedean pieces at basepoint $\alpha$, a statement about one explicit function of $\theta$ for each $\alpha$; and the Poisson form shows a zero to the right of $\alpha$ as a Cauchy kernel of negative width, whose signature Section~\ref{ch:offline} calibrates exactly.

Step 8 carries the GGC property from a fixed $\alpha$ to every $\alpha>\frac12$. Weak limits alone cannot do this (Section~\ref{ch:weak}), because the limit of the Cauchy flow is positive whether or not RH holds. The three mechanisms that could do this (a positive approximating family, an operator whose spectrum is the set of ordinates, a direct inequality for the variogram pairing) are discussed in Sections~\ref{route:sec:missing} and \ref{route:sec:plan}. Once the first $N$ zeros are verified, RH is equivalent to the GGC property of the tail factor left when their factors are removed from the clock \cite[Corollary 18.11]{P26}.

Spin chains give finite approximations of another kind: the partition functions of Knauf's chain are Dirichlet polynomials that converge to $\zeta(\beta-1)/\zeta(\beta)$ for $\re\beta>2$, and by Remark~\ref{fm:rem:neutral} the convergence of their corrected form on a strip in $1<\re\beta<2$ would carry their zeros to the zeros $1+\rho$ of $\zeta(\beta-1)$ in it, a consequence that does not decide RH, while its convergence on $\re\beta>\frac12$ without the point $2$, where the zeros of $\zeta(\beta)$ are poles of the limit, would imply it (Section~\ref{ch:fmchain}).

\begin{remark}[a Gaussian process: RH equivalent and circular]\label{rem:ch25:gp}
RH holds if and only if there is a centred Gaussian process $X$ on $[0,1]$ with continuous covariance such that $\E\exp(-\lambda\int_0^1X_t^2dt)=(\xi(\frac12)/\xi(\frac12+\sqrt{2\lambda}))^{1/2}$ for $\lambda>0$ (an unrefereed note of Polson, November 2026; proved). Under RH take the covariance $\sum_n\gamma_n^{-2}e_n(s)e_n(t)$ with $e_n=\sqrt2\sin n\pi t$, whose eigenvalues are $\gamma_n^{-2}$, and use the Hadamard product. Conversely such an $X$ gives $\xi(\frac12+\sqrt{2\lambda})/\xi(\frac12)=\prod(1+2\lambda\mu_n)$ with $\mu_n\ge0$, whose zeros in $\lambda$ are real and negative, so every ordinate is real. The covariance is built from the zeros, so the statement restates RH and does not approach it. Under RH it has an empirical process form, $T_n=\sum_k(k\pi/\gamma_k)^2\langle\alpha_n,e_k\rangle^2\Rightarrow\sum_kZ_k^2/\gamma_k^2$ for the uniform empirical process $\alpha_n$, a Fourier weighted Cram\'er von Mises statistic with weights of order $\frac14(\log k)^2$, defined only under RH.
\end{remark}

\emph{Where the self decomposable law sits} (open): the note takes the reciprocal $\Xi$ law to be self decomposable unconditionally, and a GGC under RH; whether it lies unconditionally in the smaller Urbanik and Sato classes $L_m$ inside the self decomposable laws, and what that says about the rungs of Theorem~\ref{thm:ch12:cm}, is not known.

